\documentclass[12pt]{article}
\usepackage{verbatim}
\usepackage[dvips]{epsfig}
\usepackage{color}
\usepackage{url}
\usepackage[colorlinks=true]{hyperref}

\begin{document}

\section*{Neurospaces: Documentation}

{\bf Related Documentation:}
% start: content-framework-tagreplaceitems related-do-nothing
% end: content-framework-tagreplaceitems related-do-nothing

\section*{Content framework cron jobs}

The Neurospaces content framework supports the automated construction
of websites using cron jobs (see also
\href{http://www.genesis-sim.org/userdocs/neurospaces-cron/neurospaces-cron.html}{the
  Genesis-3 cron job manual} and
\href{http://www.genesis-sim.org/userdocs/userdocs-cron/userdocs-cron.html}{the
  Genesis-3 Documentation System cron job manual}).  Cron jobs can be
associated with specific packages that have previously been
instantiated with the command {\tt neurospaces\_new\_component
  --content-framework}.


\subsection*{To create a new cron job}

\begin{enumerate}
\item Create a new user specific for the cron job with the {\tt
    adduser} unix shell command.

  Note that the repository of new packages can be served from their
  own account.  If that is the case, make sure that the cron job does
  not use this same account as that may lead to monotone database lock
  problems.

\item If the cron job is going to reinstall any of the packages, give
  the new user passwordless {\tt sudo} access: {\tt echo 'joeuser
    ALL=(ALL) NOPASSWD:ALL' >> /etc/sudoers}

\item Put the new user into the {\tt content-framework} group to make
  sure it has write access to the html output directory.
\item Create a suitable configuration file for the cron job.  Sample
  contents:
\begin{verbatim}
---
MAILTO: hugo.cornelis@gmail.com, allan.coop@gmail.com
MAILFROM: hfof_cron
HTMLDIR: /var/www/html
OUTPUTDIR: /tmp
URL: http://91.183.94.6/html/hfof
LOGROTATE: 1000
\end{verbatim}
\item Change your role to the new user.
\item {\tt neurospaces\_create\_directories <package-name>}
\item {\tt neurospaces\_pull <package-name>} {\it note: this step should
    become optional in the future}
\item {\tt neurospaces\_upgrade <package-name>}

  This gives you the {\tt <package-name>\_cron} command.  You can test it
  in isolation by running it from the command line.
\item Add {\tt <package-name>\_cron} to the {\tt /etc/crontab} file, for
  instance for execution every hour:

\begin{verbatim}
1 *     * * *   <user-name>  ( cd /home/<user-name>/neurospaces_project/<package-name>/source/snapshots/0 && /usr/local/bin/<package-name>_cron --config /home/<username>/<package-name>_cron.yml )
\end{verbatim}
\end{enumerate}


%\subsubsection*{Figures}

%Figures can easily be added with the following code:

%\begin{verbatim}
%\begin{figure}[h]
%   \centering
%   \includegraphics[scale=0.6]{figures/dummyfig.eps}
%   \caption{{\bf A Dummy Figure:} Example of Latex code to incorporate
%      a figure into documentation.}
%   \label{fig:df-1}
%\end{figure}
%\end{verbatim}

%\begin{figure}[h]
%  \centering
%   \includegraphics[scale=0.5]{figures/dummyfig.eps}
%\caption{{\bf A Dummy Figure:} Example of \LaTeX\,\,\,code to incorporate a figure into documentation.}
%  \label{fig:df-2}
%\end{figure}

\bibliographystyle{plain}
\bibliography{../tex/bib/g3-refs.bib}

\end{document}
